\documentclass[aspectratio=169]{beamer}
\usetheme{metropolis}
\metroset{block=fill}
\usefonttheme{professionalfonts}
\usepackage{appendixnumberbeamer}
% metropolis sets its title page in a box of the full paper height, which
% LaTeX reports as an overfull \vbox; this lets the box shrink to the text
% height (the title block sits 11 pt lower, nothing else changes).
\makeatletter
\expandafter\let\expandafter\metropolistitlepage\csname beamer@@tmpl@title page\endcsname
\expandafter\def\csname beamer@@tmpl@title page\endcsname{\vbox to\textheight{\metropolistitlepage\vss}}
\makeatother

% Speaker notes. deck.pdf is the audience copy; build the presenter copy, with
% each slide's notes on a second screen to the right, as
%   lualatex -jobname=deck-notes '\def\withnotes{}\documentclass[aspectratio=169]{beamer}
\usetheme{metropolis}
\metroset{block=fill}
\usefonttheme{professionalfonts}
\usepackage{appendixnumberbeamer}
% metropolis sets its title page in a box of the full paper height, which
% LaTeX reports as an overfull \vbox; this lets the box shrink to the text
% height (the title block sits 11 pt lower, nothing else changes).
\makeatletter
\expandafter\let\expandafter\metropolistitlepage\csname beamer@@tmpl@title page\endcsname
\expandafter\def\csname beamer@@tmpl@title page\endcsname{\vbox to\textheight{\metropolistitlepage\vss}}
\makeatother

% Speaker notes. deck.pdf is the audience copy; build the presenter copy, with
% each slide's notes on a second screen to the right, as
%   lualatex -jobname=deck-notes '\def\withnotes{}\documentclass[aspectratio=169]{beamer}
\usetheme{metropolis}
\metroset{block=fill}
\usefonttheme{professionalfonts}
\usepackage{appendixnumberbeamer}
% metropolis sets its title page in a box of the full paper height, which
% LaTeX reports as an overfull \vbox; this lets the box shrink to the text
% height (the title block sits 11 pt lower, nothing else changes).
\makeatletter
\expandafter\let\expandafter\metropolistitlepage\csname beamer@@tmpl@title page\endcsname
\expandafter\def\csname beamer@@tmpl@title page\endcsname{\vbox to\textheight{\metropolistitlepage\vss}}
\makeatother

% Speaker notes. deck.pdf is the audience copy; build the presenter copy, with
% each slide's notes on a second screen to the right, as
%   lualatex -jobname=deck-notes '\def\withnotes{}\documentclass[aspectratio=169]{beamer}
\usetheme{metropolis}
\metroset{block=fill}
\usefonttheme{professionalfonts}
\usepackage{appendixnumberbeamer}
% metropolis sets its title page in a box of the full paper height, which
% LaTeX reports as an overfull \vbox; this lets the box shrink to the text
% height (the title block sits 11 pt lower, nothing else changes).
\makeatletter
\expandafter\let\expandafter\metropolistitlepage\csname beamer@@tmpl@title page\endcsname
\expandafter\def\csname beamer@@tmpl@title page\endcsname{\vbox to\textheight{\metropolistitlepage\vss}}
\makeatother

% Speaker notes. deck.pdf is the audience copy; build the presenter copy, with
% each slide's notes on a second screen to the right, as
%   lualatex -jobname=deck-notes '\def\withnotes{}\input{deck}'
\ifdefined\withnotes
  \usepackage{pgfpages}
  \setbeameroption{show notes on second screen=right}
\fi
\setbeamertemplate{note page}{\vspace{4mm}\small\insertnote}

% Catppuccin Latte and the theme colours.
\usepackage{xcolor}
\input{colors.tex}
% \usepreset{mpp}
\setbeamerfont{button}{size=\scriptsize}
\setbeamerfont{page number in head/foot}{size=\scriptsize}
\metroset{numbering=counter}
% A navigation strip of the written work goes left of the frame number on
% frames about one step, e.g. in a group around them:
%   \setbeamertemplate{frame footer}{\resizebox{!}{4mm}{\input{assets/strip.tex}}}

% Build with LuaLaTeX: metropolis then sets its Fira Sans (no fontenc needed).
\usepackage[USenglish]{babel}
\usepackage{amsmath,amssymb,mathtools}
\usepackage{siunitx}
\sisetup{separate-uncertainty,locale=US,mode=text}
\usepackage{graphicx}
\usepackage{booktabs}
\usepackage{csquotes}
\usepackage{tikz}
\usetikzlibrary{arrows.meta,positioning}

\usepackage[backend=biber,sorting=none,backref=false]{biblatex}
\renewcommand*{\bibfont}{\footnotesize}
% Point this at the shared bibliography library; no deck-local copies of sources.
\addbibresource{../bib/library.bib}

\hypersetup{colorlinks=true,linkcolor=DeckAccent,urlcolor=CtpFlamingo,citecolor=CtpTeal}

\title{Title that states the result}
\subtitle{Meeting or course}
\date{1 January 2026}
\author{Author}
% Logos and institute are personal settings, e.g. a white logo on the title page:
% \institute{\null\hfill\includegraphics[height=8mm]{logo.pdf}}

\begin{document}

{
\setbeamercolor{background canvas}{bg=DeckMain}
\setbeamercolor{normal text}{fg=CtpBase}
\maketitle
}

% The question the talk answers, alone on one slide.
\begin{frame}{}
  \centering\vfill\Large
  \textbf{The question this talk answers?}
  \vfill
  \note{Say why the audience should care in two sentences.}
\end{frame}

\section{Motivation}

% The frame title states the slide's one claim; the body supports it.
\begin{frame}[label=claim]{The title states the claim}
  \begin{columns}
    \begin{column}{0.5\textwidth}
      \begin{itemize}
        \item One idea per slide
        \item Words the speaker will say, not the script
      \end{itemize}
      \pause
      % A key number or definition in a filled box; terms commented in place.
      \begin{block}{The relation behind the claim}
        \[ \varepsilon = \frac{\overbrace{N_\text{pass}}^{\text{selected}}}{\underbrace{N_\text{all}}_{\text{generated}}} \]
      \end{block}
      \textcolor{CtpRed}{The point the audience must remember.}
    \end{column}
    \begin{column}{0.45\textwidth}
      % \includegraphics[width=\textwidth]{figure.pdf}
      \scriptsize Source: \cite{key} or the study that made the figure.
    \end{column}
  \end{columns}
  \vfill
  \hyperlink{backup-details}{\beamergotobutton{Details}}
  \note{What to say, what to point at, and the time mark: 2 min.}
\end{frame}

\section{Conclusion}

% The take-away, one sentence, before the references.
\begin{frame}
  \begin{block}{Take-away}
    \Large One sentence the audience should repeat afterwards.
  \end{block}
  \note{Pause here; this is the slide left up during questions if there is no
    questions slide.}
\end{frame}

\begin{frame}[allowframebreaks]
  \printbibliography
\end{frame}

{
\setbeamercolor{background canvas}{bg=DeckMain}
\begin{frame}[plain]
  \centering\vfill\huge
  \textcolor{CtpBase}{\textbf{Questions?}}
  \vfill
\end{frame}
}

\appendix

% Backup slides answer the questions you expect; each links back.
\begin{frame}[label=backup-details]{Backup: the details behind the claim}
  Material for an expected question.
  \vfill\hfill\hyperlink{claim}{\beamerreturnbutton{Back}}
\end{frame}

% A blank last page to draw on during questions.
\begin{frame}[plain,noframenumbering]
  \null
\end{frame}

\end{document}
'
\ifdefined\withnotes
  \usepackage{pgfpages}
  \setbeameroption{show notes on second screen=right}
\fi
\setbeamertemplate{note page}{\vspace{4mm}\small\insertnote}

% Catppuccin Latte and the theme colours.
\usepackage{xcolor}
% Colorscheme
\definecolor{CtpRosewater}{RGB}{220,138,120}
\definecolor{CtpFlamingo}{RGB}{221,120,120}
\definecolor{CtpPink}{RGB}{234,118,203}
\definecolor{CtpMauve}{RGB}{136,57,239}
\definecolor{CtpRed}{RGB}{210,15,57}
\definecolor{CtpMaroon}{RGB}{230,69,83}
\definecolor{CtpPeach}{RGB}{254,100,11}
\definecolor{CtpYellow}{RGB}{223,142,29}
\definecolor{CtpGreen}{RGB}{64,160,43}
\definecolor{CtpTeal}{RGB}{23,146,153}
\definecolor{CtpSky}{RGB}{4,165,229}
\definecolor{CtpSapphire}{RGB}{32,159,181}
\definecolor{CtpBlue}{RGB}{30,102,245}
\definecolor{CtpLavender}{RGB}{114,135,253}
\definecolor{CtpText}{RGB}{76,79,105}
\definecolor{CtpSubtext1}{RGB}{92,95,119}
\definecolor{CtpSubtext0}{RGB}{108,111,133}
\definecolor{CtpOverlay2}{RGB}{124,127,147}
\definecolor{CtpOverlay1}{RGB}{140,143,161}
\definecolor{CtpOverlay0}{RGB}{156,160,176}
\definecolor{CtpSurface2}{RGB}{172,176,190}
\definecolor{CtpSurface1}{RGB}{188,192,204}
\definecolor{CtpSurface0}{RGB}{204,208,218}
\definecolor{CtpBase}{RGB}{239,241,245}
\definecolor{CtpMantle}{RGB}{230,233,239}
\definecolor{CtpCrust}{RGB}{220,224,232}

% Institution presets: a main colour (title page, title bar, progress bar,
% buttons) and an accent (internal links, which include metropolis's section
% titles). TUM is the default; choose
% another after \input{colors.tex}, e.g. \usepreset{mpp}.
\newcommand{\definepreset}[3]{\expandafter\def\csname preset:#1\endcsname{\colorlet{DeckMain}{#2}\colorlet{DeckAccent}{#3}}}
\definepreset{tum}{CtpBlue}{CtpSky}
\definepreset{mpp}{CtpTeal}{CtpGreen}
\newcommand{\usepreset}[1]{\csname preset:#1\endcsname}
\usepreset{tum}

\setbeamercolor{normal text}{fg=black}
\setbeamercolor{alerted text}{fg=CtpYellow}
\setbeamercolor{frametitle}{fg=CtpBase, bg=DeckMain}
\setbeamercolor{progress bar}{fg=DeckMain, bg=CtpSurface0}
\setbeamercolor{button}{fg=white, bg=DeckMain}
\setbeamercolor{section title}{fg=CtpText}
\setbeamerfont{caption}{size=\tiny} % Use `\tiny`, `\scriptsize`, etc., for even smaller text
\setbeamercolor{caption}{fg=CtpText} % Change color to gray


% \usepreset{mpp}
\setbeamerfont{button}{size=\scriptsize}
\setbeamerfont{page number in head/foot}{size=\scriptsize}
\metroset{numbering=counter}
% A navigation strip of the written work goes left of the frame number on
% frames about one step, e.g. in a group around them:
%   \setbeamertemplate{frame footer}{\resizebox{!}{4mm}{\input{assets/strip.tex}}}

% Build with LuaLaTeX: metropolis then sets its Fira Sans (no fontenc needed).
\usepackage[USenglish]{babel}
\usepackage{amsmath,amssymb,mathtools}
\usepackage{siunitx}
\sisetup{separate-uncertainty,locale=US,mode=text}
\usepackage{graphicx}
\usepackage{booktabs}
\usepackage{csquotes}
\usepackage{tikz}
\usetikzlibrary{arrows.meta,positioning}

\usepackage[backend=biber,sorting=none,backref=false]{biblatex}
\renewcommand*{\bibfont}{\footnotesize}
% Point this at the shared bibliography library; no deck-local copies of sources.
\addbibresource{../bib/library.bib}

\hypersetup{colorlinks=true,linkcolor=DeckAccent,urlcolor=CtpFlamingo,citecolor=CtpTeal}

\title{Title that states the result}
\subtitle{Meeting or course}
\date{1 January 2026}
\author{Author}
% Logos and institute are personal settings, e.g. a white logo on the title page:
% \institute{\null\hfill\includegraphics[height=8mm]{logo.pdf}}

\begin{document}

{
\setbeamercolor{background canvas}{bg=DeckMain}
\setbeamercolor{normal text}{fg=CtpBase}
\maketitle
}

% The question the talk answers, alone on one slide.
\begin{frame}{}
  \centering\vfill\Large
  \textbf{The question this talk answers?}
  \vfill
  \note{Say why the audience should care in two sentences.}
\end{frame}

\section{Motivation}

% The frame title states the slide's one claim; the body supports it.
\begin{frame}[label=claim]{The title states the claim}
  \begin{columns}
    \begin{column}{0.5\textwidth}
      \begin{itemize}
        \item One idea per slide
        \item Words the speaker will say, not the script
      \end{itemize}
      \pause
      % A key number or definition in a filled box; terms commented in place.
      \begin{block}{The relation behind the claim}
        \[ \varepsilon = \frac{\overbrace{N_\text{pass}}^{\text{selected}}}{\underbrace{N_\text{all}}_{\text{generated}}} \]
      \end{block}
      \textcolor{CtpRed}{The point the audience must remember.}
    \end{column}
    \begin{column}{0.45\textwidth}
      % \includegraphics[width=\textwidth]{figure.pdf}
      \scriptsize Source: \cite{key} or the study that made the figure.
    \end{column}
  \end{columns}
  \vfill
  \hyperlink{backup-details}{\beamergotobutton{Details}}
  \note{What to say, what to point at, and the time mark: 2 min.}
\end{frame}

\section{Conclusion}

% The take-away, one sentence, before the references.
\begin{frame}
  \begin{block}{Take-away}
    \Large One sentence the audience should repeat afterwards.
  \end{block}
  \note{Pause here; this is the slide left up during questions if there is no
    questions slide.}
\end{frame}

\begin{frame}[allowframebreaks]
  \printbibliography
\end{frame}

{
\setbeamercolor{background canvas}{bg=DeckMain}
\begin{frame}[plain]
  \centering\vfill\huge
  \textcolor{CtpBase}{\textbf{Questions?}}
  \vfill
\end{frame}
}

\appendix

% Backup slides answer the questions you expect; each links back.
\begin{frame}[label=backup-details]{Backup: the details behind the claim}
  Material for an expected question.
  \vfill\hfill\hyperlink{claim}{\beamerreturnbutton{Back}}
\end{frame}

% A blank last page to draw on during questions.
\begin{frame}[plain,noframenumbering]
  \null
\end{frame}

\end{document}
'
\ifdefined\withnotes
  \usepackage{pgfpages}
  \setbeameroption{show notes on second screen=right}
\fi
\setbeamertemplate{note page}{\vspace{4mm}\small\insertnote}

% Catppuccin Latte and the theme colours.
\usepackage{xcolor}
% Colorscheme
\definecolor{CtpRosewater}{RGB}{220,138,120}
\definecolor{CtpFlamingo}{RGB}{221,120,120}
\definecolor{CtpPink}{RGB}{234,118,203}
\definecolor{CtpMauve}{RGB}{136,57,239}
\definecolor{CtpRed}{RGB}{210,15,57}
\definecolor{CtpMaroon}{RGB}{230,69,83}
\definecolor{CtpPeach}{RGB}{254,100,11}
\definecolor{CtpYellow}{RGB}{223,142,29}
\definecolor{CtpGreen}{RGB}{64,160,43}
\definecolor{CtpTeal}{RGB}{23,146,153}
\definecolor{CtpSky}{RGB}{4,165,229}
\definecolor{CtpSapphire}{RGB}{32,159,181}
\definecolor{CtpBlue}{RGB}{30,102,245}
\definecolor{CtpLavender}{RGB}{114,135,253}
\definecolor{CtpText}{RGB}{76,79,105}
\definecolor{CtpSubtext1}{RGB}{92,95,119}
\definecolor{CtpSubtext0}{RGB}{108,111,133}
\definecolor{CtpOverlay2}{RGB}{124,127,147}
\definecolor{CtpOverlay1}{RGB}{140,143,161}
\definecolor{CtpOverlay0}{RGB}{156,160,176}
\definecolor{CtpSurface2}{RGB}{172,176,190}
\definecolor{CtpSurface1}{RGB}{188,192,204}
\definecolor{CtpSurface0}{RGB}{204,208,218}
\definecolor{CtpBase}{RGB}{239,241,245}
\definecolor{CtpMantle}{RGB}{230,233,239}
\definecolor{CtpCrust}{RGB}{220,224,232}

% Institution presets: a main colour (title page, title bar, progress bar,
% buttons) and an accent (internal links, which include metropolis's section
% titles). TUM is the default; choose
% another after % Colorscheme
\definecolor{CtpRosewater}{RGB}{220,138,120}
\definecolor{CtpFlamingo}{RGB}{221,120,120}
\definecolor{CtpPink}{RGB}{234,118,203}
\definecolor{CtpMauve}{RGB}{136,57,239}
\definecolor{CtpRed}{RGB}{210,15,57}
\definecolor{CtpMaroon}{RGB}{230,69,83}
\definecolor{CtpPeach}{RGB}{254,100,11}
\definecolor{CtpYellow}{RGB}{223,142,29}
\definecolor{CtpGreen}{RGB}{64,160,43}
\definecolor{CtpTeal}{RGB}{23,146,153}
\definecolor{CtpSky}{RGB}{4,165,229}
\definecolor{CtpSapphire}{RGB}{32,159,181}
\definecolor{CtpBlue}{RGB}{30,102,245}
\definecolor{CtpLavender}{RGB}{114,135,253}
\definecolor{CtpText}{RGB}{76,79,105}
\definecolor{CtpSubtext1}{RGB}{92,95,119}
\definecolor{CtpSubtext0}{RGB}{108,111,133}
\definecolor{CtpOverlay2}{RGB}{124,127,147}
\definecolor{CtpOverlay1}{RGB}{140,143,161}
\definecolor{CtpOverlay0}{RGB}{156,160,176}
\definecolor{CtpSurface2}{RGB}{172,176,190}
\definecolor{CtpSurface1}{RGB}{188,192,204}
\definecolor{CtpSurface0}{RGB}{204,208,218}
\definecolor{CtpBase}{RGB}{239,241,245}
\definecolor{CtpMantle}{RGB}{230,233,239}
\definecolor{CtpCrust}{RGB}{220,224,232}

% Institution presets: a main colour (title page, title bar, progress bar,
% buttons) and an accent (internal links, which include metropolis's section
% titles). TUM is the default; choose
% another after \input{colors.tex}, e.g. \usepreset{mpp}.
\newcommand{\definepreset}[3]{\expandafter\def\csname preset:#1\endcsname{\colorlet{DeckMain}{#2}\colorlet{DeckAccent}{#3}}}
\definepreset{tum}{CtpBlue}{CtpSky}
\definepreset{mpp}{CtpTeal}{CtpGreen}
\newcommand{\usepreset}[1]{\csname preset:#1\endcsname}
\usepreset{tum}

\setbeamercolor{normal text}{fg=black}
\setbeamercolor{alerted text}{fg=CtpYellow}
\setbeamercolor{frametitle}{fg=CtpBase, bg=DeckMain}
\setbeamercolor{progress bar}{fg=DeckMain, bg=CtpSurface0}
\setbeamercolor{button}{fg=white, bg=DeckMain}
\setbeamercolor{section title}{fg=CtpText}
\setbeamerfont{caption}{size=\tiny} % Use `\tiny`, `\scriptsize`, etc., for even smaller text
\setbeamercolor{caption}{fg=CtpText} % Change color to gray

, e.g. \usepreset{mpp}.
\newcommand{\definepreset}[3]{\expandafter\def\csname preset:#1\endcsname{\colorlet{DeckMain}{#2}\colorlet{DeckAccent}{#3}}}
\definepreset{tum}{CtpBlue}{CtpSky}
\definepreset{mpp}{CtpTeal}{CtpGreen}
\newcommand{\usepreset}[1]{\csname preset:#1\endcsname}
\usepreset{tum}

\setbeamercolor{normal text}{fg=black}
\setbeamercolor{alerted text}{fg=CtpYellow}
\setbeamercolor{frametitle}{fg=CtpBase, bg=DeckMain}
\setbeamercolor{progress bar}{fg=DeckMain, bg=CtpSurface0}
\setbeamercolor{button}{fg=white, bg=DeckMain}
\setbeamercolor{section title}{fg=CtpText}
\setbeamerfont{caption}{size=\tiny} % Use `\tiny`, `\scriptsize`, etc., for even smaller text
\setbeamercolor{caption}{fg=CtpText} % Change color to gray


% \usepreset{mpp}
\setbeamerfont{button}{size=\scriptsize}
\setbeamerfont{page number in head/foot}{size=\scriptsize}
\metroset{numbering=counter}
% A navigation strip of the written work goes left of the frame number on
% frames about one step, e.g. in a group around them:
%   \setbeamertemplate{frame footer}{\resizebox{!}{4mm}{\input{assets/strip.tex}}}

% Build with LuaLaTeX: metropolis then sets its Fira Sans (no fontenc needed).
\usepackage[USenglish]{babel}
\usepackage{amsmath,amssymb,mathtools}
\usepackage{siunitx}
\sisetup{separate-uncertainty,locale=US,mode=text}
\usepackage{graphicx}
\usepackage{booktabs}
\usepackage{csquotes}
\usepackage{tikz}
\usetikzlibrary{arrows.meta,positioning}

\usepackage[backend=biber,sorting=none,backref=false]{biblatex}
\renewcommand*{\bibfont}{\footnotesize}
% Point this at the shared bibliography library; no deck-local copies of sources.
\addbibresource{../bib/library.bib}

\hypersetup{colorlinks=true,linkcolor=DeckAccent,urlcolor=CtpFlamingo,citecolor=CtpTeal}

\title{Title that states the result}
\subtitle{Meeting or course}
\date{1 January 2026}
\author{Author}
% Logos and institute are personal settings, e.g. a white logo on the title page:
% \institute{\null\hfill\includegraphics[height=8mm]{logo.pdf}}

\begin{document}

{
\setbeamercolor{background canvas}{bg=DeckMain}
\setbeamercolor{normal text}{fg=CtpBase}
\maketitle
}

% The question the talk answers, alone on one slide.
\begin{frame}{}
  \centering\vfill\Large
  \textbf{The question this talk answers?}
  \vfill
  \note{Say why the audience should care in two sentences.}
\end{frame}

\section{Motivation}

% The frame title states the slide's one claim; the body supports it.
\begin{frame}[label=claim]{The title states the claim}
  \begin{columns}
    \begin{column}{0.5\textwidth}
      \begin{itemize}
        \item One idea per slide
        \item Words the speaker will say, not the script
      \end{itemize}
      \pause
      % A key number or definition in a filled box; terms commented in place.
      \begin{block}{The relation behind the claim}
        \[ \varepsilon = \frac{\overbrace{N_\text{pass}}^{\text{selected}}}{\underbrace{N_\text{all}}_{\text{generated}}} \]
      \end{block}
      \textcolor{CtpRed}{The point the audience must remember.}
    \end{column}
    \begin{column}{0.45\textwidth}
      % \includegraphics[width=\textwidth]{figure.pdf}
      \scriptsize Source: \cite{key} or the study that made the figure.
    \end{column}
  \end{columns}
  \vfill
  \hyperlink{backup-details}{\beamergotobutton{Details}}
  \note{What to say, what to point at, and the time mark: 2 min.}
\end{frame}

\section{Conclusion}

% The take-away, one sentence, before the references.
\begin{frame}
  \begin{block}{Take-away}
    \Large One sentence the audience should repeat afterwards.
  \end{block}
  \note{Pause here; this is the slide left up during questions if there is no
    questions slide.}
\end{frame}

\begin{frame}[allowframebreaks]
  \printbibliography
\end{frame}

{
\setbeamercolor{background canvas}{bg=DeckMain}
\begin{frame}[plain]
  \centering\vfill\huge
  \textcolor{CtpBase}{\textbf{Questions?}}
  \vfill
\end{frame}
}

\appendix

% Backup slides answer the questions you expect; each links back.
\begin{frame}[label=backup-details]{Backup: the details behind the claim}
  Material for an expected question.
  \vfill\hfill\hyperlink{claim}{\beamerreturnbutton{Back}}
\end{frame}

% A blank last page to draw on during questions.
\begin{frame}[plain,noframenumbering]
  \null
\end{frame}

\end{document}
'
\ifdefined\withnotes
  \usepackage{pgfpages}
  \setbeameroption{show notes on second screen=right}
\fi
\setbeamertemplate{note page}{\vspace{4mm}\small\insertnote}

% Catppuccin Latte and the theme colours.
\usepackage{xcolor}
% Colorscheme
\definecolor{CtpRosewater}{RGB}{220,138,120}
\definecolor{CtpFlamingo}{RGB}{221,120,120}
\definecolor{CtpPink}{RGB}{234,118,203}
\definecolor{CtpMauve}{RGB}{136,57,239}
\definecolor{CtpRed}{RGB}{210,15,57}
\definecolor{CtpMaroon}{RGB}{230,69,83}
\definecolor{CtpPeach}{RGB}{254,100,11}
\definecolor{CtpYellow}{RGB}{223,142,29}
\definecolor{CtpGreen}{RGB}{64,160,43}
\definecolor{CtpTeal}{RGB}{23,146,153}
\definecolor{CtpSky}{RGB}{4,165,229}
\definecolor{CtpSapphire}{RGB}{32,159,181}
\definecolor{CtpBlue}{RGB}{30,102,245}
\definecolor{CtpLavender}{RGB}{114,135,253}
\definecolor{CtpText}{RGB}{76,79,105}
\definecolor{CtpSubtext1}{RGB}{92,95,119}
\definecolor{CtpSubtext0}{RGB}{108,111,133}
\definecolor{CtpOverlay2}{RGB}{124,127,147}
\definecolor{CtpOverlay1}{RGB}{140,143,161}
\definecolor{CtpOverlay0}{RGB}{156,160,176}
\definecolor{CtpSurface2}{RGB}{172,176,190}
\definecolor{CtpSurface1}{RGB}{188,192,204}
\definecolor{CtpSurface0}{RGB}{204,208,218}
\definecolor{CtpBase}{RGB}{239,241,245}
\definecolor{CtpMantle}{RGB}{230,233,239}
\definecolor{CtpCrust}{RGB}{220,224,232}

% Institution presets: a main colour (title page, title bar, progress bar,
% buttons) and an accent (internal links, which include metropolis's section
% titles). TUM is the default; choose
% another after % Colorscheme
\definecolor{CtpRosewater}{RGB}{220,138,120}
\definecolor{CtpFlamingo}{RGB}{221,120,120}
\definecolor{CtpPink}{RGB}{234,118,203}
\definecolor{CtpMauve}{RGB}{136,57,239}
\definecolor{CtpRed}{RGB}{210,15,57}
\definecolor{CtpMaroon}{RGB}{230,69,83}
\definecolor{CtpPeach}{RGB}{254,100,11}
\definecolor{CtpYellow}{RGB}{223,142,29}
\definecolor{CtpGreen}{RGB}{64,160,43}
\definecolor{CtpTeal}{RGB}{23,146,153}
\definecolor{CtpSky}{RGB}{4,165,229}
\definecolor{CtpSapphire}{RGB}{32,159,181}
\definecolor{CtpBlue}{RGB}{30,102,245}
\definecolor{CtpLavender}{RGB}{114,135,253}
\definecolor{CtpText}{RGB}{76,79,105}
\definecolor{CtpSubtext1}{RGB}{92,95,119}
\definecolor{CtpSubtext0}{RGB}{108,111,133}
\definecolor{CtpOverlay2}{RGB}{124,127,147}
\definecolor{CtpOverlay1}{RGB}{140,143,161}
\definecolor{CtpOverlay0}{RGB}{156,160,176}
\definecolor{CtpSurface2}{RGB}{172,176,190}
\definecolor{CtpSurface1}{RGB}{188,192,204}
\definecolor{CtpSurface0}{RGB}{204,208,218}
\definecolor{CtpBase}{RGB}{239,241,245}
\definecolor{CtpMantle}{RGB}{230,233,239}
\definecolor{CtpCrust}{RGB}{220,224,232}

% Institution presets: a main colour (title page, title bar, progress bar,
% buttons) and an accent (internal links, which include metropolis's section
% titles). TUM is the default; choose
% another after % Colorscheme
\definecolor{CtpRosewater}{RGB}{220,138,120}
\definecolor{CtpFlamingo}{RGB}{221,120,120}
\definecolor{CtpPink}{RGB}{234,118,203}
\definecolor{CtpMauve}{RGB}{136,57,239}
\definecolor{CtpRed}{RGB}{210,15,57}
\definecolor{CtpMaroon}{RGB}{230,69,83}
\definecolor{CtpPeach}{RGB}{254,100,11}
\definecolor{CtpYellow}{RGB}{223,142,29}
\definecolor{CtpGreen}{RGB}{64,160,43}
\definecolor{CtpTeal}{RGB}{23,146,153}
\definecolor{CtpSky}{RGB}{4,165,229}
\definecolor{CtpSapphire}{RGB}{32,159,181}
\definecolor{CtpBlue}{RGB}{30,102,245}
\definecolor{CtpLavender}{RGB}{114,135,253}
\definecolor{CtpText}{RGB}{76,79,105}
\definecolor{CtpSubtext1}{RGB}{92,95,119}
\definecolor{CtpSubtext0}{RGB}{108,111,133}
\definecolor{CtpOverlay2}{RGB}{124,127,147}
\definecolor{CtpOverlay1}{RGB}{140,143,161}
\definecolor{CtpOverlay0}{RGB}{156,160,176}
\definecolor{CtpSurface2}{RGB}{172,176,190}
\definecolor{CtpSurface1}{RGB}{188,192,204}
\definecolor{CtpSurface0}{RGB}{204,208,218}
\definecolor{CtpBase}{RGB}{239,241,245}
\definecolor{CtpMantle}{RGB}{230,233,239}
\definecolor{CtpCrust}{RGB}{220,224,232}

% Institution presets: a main colour (title page, title bar, progress bar,
% buttons) and an accent (internal links, which include metropolis's section
% titles). TUM is the default; choose
% another after \input{colors.tex}, e.g. \usepreset{mpp}.
\newcommand{\definepreset}[3]{\expandafter\def\csname preset:#1\endcsname{\colorlet{DeckMain}{#2}\colorlet{DeckAccent}{#3}}}
\definepreset{tum}{CtpBlue}{CtpSky}
\definepreset{mpp}{CtpTeal}{CtpGreen}
\newcommand{\usepreset}[1]{\csname preset:#1\endcsname}
\usepreset{tum}

\setbeamercolor{normal text}{fg=black}
\setbeamercolor{alerted text}{fg=CtpYellow}
\setbeamercolor{frametitle}{fg=CtpBase, bg=DeckMain}
\setbeamercolor{progress bar}{fg=DeckMain, bg=CtpSurface0}
\setbeamercolor{button}{fg=white, bg=DeckMain}
\setbeamercolor{section title}{fg=CtpText}
\setbeamerfont{caption}{size=\tiny} % Use `\tiny`, `\scriptsize`, etc., for even smaller text
\setbeamercolor{caption}{fg=CtpText} % Change color to gray

, e.g. \usepreset{mpp}.
\newcommand{\definepreset}[3]{\expandafter\def\csname preset:#1\endcsname{\colorlet{DeckMain}{#2}\colorlet{DeckAccent}{#3}}}
\definepreset{tum}{CtpBlue}{CtpSky}
\definepreset{mpp}{CtpTeal}{CtpGreen}
\newcommand{\usepreset}[1]{\csname preset:#1\endcsname}
\usepreset{tum}

\setbeamercolor{normal text}{fg=black}
\setbeamercolor{alerted text}{fg=CtpYellow}
\setbeamercolor{frametitle}{fg=CtpBase, bg=DeckMain}
\setbeamercolor{progress bar}{fg=DeckMain, bg=CtpSurface0}
\setbeamercolor{button}{fg=white, bg=DeckMain}
\setbeamercolor{section title}{fg=CtpText}
\setbeamerfont{caption}{size=\tiny} % Use `\tiny`, `\scriptsize`, etc., for even smaller text
\setbeamercolor{caption}{fg=CtpText} % Change color to gray

, e.g. \usepreset{mpp}.
\newcommand{\definepreset}[3]{\expandafter\def\csname preset:#1\endcsname{\colorlet{DeckMain}{#2}\colorlet{DeckAccent}{#3}}}
\definepreset{tum}{CtpBlue}{CtpSky}
\definepreset{mpp}{CtpTeal}{CtpGreen}
\newcommand{\usepreset}[1]{\csname preset:#1\endcsname}
\usepreset{tum}

\setbeamercolor{normal text}{fg=black}
\setbeamercolor{alerted text}{fg=CtpYellow}
\setbeamercolor{frametitle}{fg=CtpBase, bg=DeckMain}
\setbeamercolor{progress bar}{fg=DeckMain, bg=CtpSurface0}
\setbeamercolor{button}{fg=white, bg=DeckMain}
\setbeamercolor{section title}{fg=CtpText}
\setbeamerfont{caption}{size=\tiny} % Use `\tiny`, `\scriptsize`, etc., for even smaller text
\setbeamercolor{caption}{fg=CtpText} % Change color to gray


% \usepreset{mpp}
\setbeamerfont{button}{size=\scriptsize}
\setbeamerfont{page number in head/foot}{size=\scriptsize}
\metroset{numbering=counter}
% A navigation strip of the written work goes left of the frame number on
% frames about one step, e.g. in a group around them:
%   \setbeamertemplate{frame footer}{\resizebox{!}{4mm}{\input{assets/strip.tex}}}

% Build with LuaLaTeX: metropolis then sets its Fira Sans (no fontenc needed).
\usepackage[USenglish]{babel}
\usepackage{amsmath,amssymb,mathtools}
\usepackage{siunitx}
\sisetup{separate-uncertainty,locale=US,mode=text}
\usepackage{graphicx}
\usepackage{booktabs}
\usepackage{csquotes}
\usepackage{tikz}
\usetikzlibrary{arrows.meta,positioning}

\usepackage[backend=biber,sorting=none,backref=false]{biblatex}
\renewcommand*{\bibfont}{\footnotesize}
% Point this at the shared bibliography library; no deck-local copies of sources.
\addbibresource{../bib/library.bib}

\hypersetup{colorlinks=true,linkcolor=DeckAccent,urlcolor=CtpFlamingo,citecolor=CtpTeal}

\title{Title that states the result}
\subtitle{Meeting or course}
\date{1 January 2026}
\author{Author}
% Logos and institute are personal settings, e.g. a white logo on the title page:
% \institute{\null\hfill\includegraphics[height=8mm]{logo.pdf}}

\begin{document}

{
\setbeamercolor{background canvas}{bg=DeckMain}
\setbeamercolor{normal text}{fg=CtpBase}
\maketitle
}

% The question the talk answers, alone on one slide.
\begin{frame}{}
  \centering\vfill\Large
  \textbf{The question this talk answers?}
  \vfill
  \note{Say why the audience should care in two sentences.}
\end{frame}

\section{Motivation}

% The frame title states the slide's one claim; the body supports it.
\begin{frame}[label=claim]{The title states the claim}
  \begin{columns}
    \begin{column}{0.5\textwidth}
      \begin{itemize}
        \item One idea per slide
        \item Words the speaker will say, not the script
      \end{itemize}
      \pause
      % A key number or definition in a filled box; terms commented in place.
      \begin{block}{The relation behind the claim}
        \[ \varepsilon = \frac{\overbrace{N_\text{pass}}^{\text{selected}}}{\underbrace{N_\text{all}}_{\text{generated}}} \]
      \end{block}
      \textcolor{CtpRed}{The point the audience must remember.}
    \end{column}
    \begin{column}{0.45\textwidth}
      % \includegraphics[width=\textwidth]{figure.pdf}
      \scriptsize Source: \cite{key} or the study that made the figure.
    \end{column}
  \end{columns}
  \vfill
  \hyperlink{backup-details}{\beamergotobutton{Details}}
  \note{What to say, what to point at, and the time mark: 2 min.}
\end{frame}

\section{Conclusion}

% The take-away, one sentence, before the references.
\begin{frame}
  \begin{block}{Take-away}
    \Large One sentence the audience should repeat afterwards.
  \end{block}
  \note{Pause here; this is the slide left up during questions if there is no
    questions slide.}
\end{frame}

\begin{frame}[allowframebreaks]
  \printbibliography
\end{frame}

{
\setbeamercolor{background canvas}{bg=DeckMain}
\begin{frame}[plain]
  \centering\vfill\huge
  \textcolor{CtpBase}{\textbf{Questions?}}
  \vfill
\end{frame}
}

\appendix

% Backup slides answer the questions you expect; each links back.
\begin{frame}[label=backup-details]{Backup: the details behind the claim}
  Material for an expected question.
  \vfill\hfill\hyperlink{claim}{\beamerreturnbutton{Back}}
\end{frame}

% A blank last page to draw on during questions.
\begin{frame}[plain,noframenumbering]
  \null
\end{frame}

\end{document}
